This paper has explored the methodological synergies between structural econometrics and reinforcement learning. We demonstrated that reinforcement learning algorithms can be viewed as generalizations of dynamic programming methods central to structural economics. RL can be used to solve static and dynamic games, static and dynamic decision making problems.

The cross-fertilization between these disciplines offers substantial benefits in both directions. Structural econometrics gains computational tools to solve, high-dimensional models previously considered intractable, while reinforcement learning benefits from economic theory's structural assumptions. One of the reasons why AlphaZero did so well was because it had a perfect model. 

An approximate and reasonably accurate model can significantly reduce the search space for learning. The success of RL in games and robotics is attributable to their existing perfect models of games and near perfect physics simulators. In economics and social sciences, we cannot expect to build perfect simulators but we can capture and exploit some empirical regularities (e.g. demand, budget constraints, loss aversion) and incorporate them to improve real-time learning algorithms. 

Methodological challenges will persist. RL is brittle, requires extensive hyperparameter tuning, lacks standardized implementations, is computationally demanding, and convergence is not guaranteed in many settings. Nevertheless, it offers a potentially powerful toolkit for dynamic and strategic decision-making.